\chapter{第五次反“围剿”：中共绝境求生}

\section{中共的反“围剿”准备和作战方针}


在国民党方面积极部署第五次“围剿”，准备向苏区发动新的进攻时，中共方面也在全力准备反“围剿”战争。对于国民党军新一轮进攻的严重性，应该说，中共方面有着充分的估计。1933年7月24日，中共中央就展开新一轮反“围剿”作出决议，表达了粉碎“围剿”的信心，认为第四次“围剿”被粉碎后，造成了国内阶级力量对比新的有利变动，国民党及其他各派军阀在政治上与军事上不断削弱，国民党统治急剧崩溃，“五次‘围剿’的粉碎，将使我们完全的可能实现中国革命一省或数省的首先胜利”。通观决议，其中确有当时习见的高调成分，但对于年轻的中共党人而言，这样的高调常常只是不愿示弱的一种表现，如果以为中共就是以此来确定其反“围剿”的战略方针，则未免低估了其能力和水准。就在这一决议中，中共中央在乐观预测总体形势同时，也严肃指出：“阶级敌人对于苏维埃和红军的新的‘围剿’，是在疯狂的准备着，五次‘围剿’是更加剧烈与残酷的阶级决战”，“要胜利的粉碎这次‘围剿’，要争取民族危机与经济浩劫中的苏维埃的出路，那我们必须紧张我们一切的努力，动员尽可能的广大的群众去参加革命斗争与革命战争。”\footnote{《中共中央委员会关于帝国主义国民党五次“围剿”与我们党的任务的决议》，《中共中央文件选集》第10册，第275～276页。}7月28日，共产国际代表埃韦特在内部报告中说得更直接：“我们没有理由悲观，但我认为，近期我们摆脱中央苏区所处困境的希望不大。”\footnote{《埃韦特给皮亚特尼茨基的第六号报告（1933年7月28日）》，《共产国际、联共（布）与中国革命档案资料丛书》第13册，第459页。}


与中共中央发出决议同日，博古在中央一级党的活动分子会议上作报告，和中共中央文件一样，博古在估计“我们有着一切的条件取得这次战争中的完全胜利，实现中国革命的一省或数省的胜利与确立独立自由的苏维埃中国”的同时，又强调指出：

\begin{quote}

在这个剧烈的战斗中间我们还有许多困难——发展的困难，假若我们不能继续我们一切的努力，紧急的布尔什维克的动员尽可能的广大群众来为着争取苏维埃的出路与粉碎敌人的五次“围剿”而斗争，为克服我们发展中的困难而斗争，那我们便不能胜利。\footnote{博古：《为粉碎敌人的五次“围剿”与争取独立自由的苏维埃中国而斗争》，《秦邦宪（博古）文集》，中共党史出版社，2007，第216页。}
\end{quote}


随后，以苏维埃中央政府名义发出的致世界劳苦民众宣言更明确宣示：“现在，这个年轻的中华苏维埃共和国却正处在一个严重的危险关头了！在这个紧急关头，我们向你们请求：帮助我们反对那些要使我们再过黑暗的非人生活，并要屠杀我们的刽子手呵。”\footnote{《中华苏维埃共和国中央政府告全世界工农劳苦民众宣言（1933年9月6日）》，竹内实编《毛泽东集》第4册，日本苍苍社，1983年第2版，第27页。}虽然，人们也许会指责在日本占领东三省后，中共中央没有及时将斗争矛头指向外国侵略者，但客观而言，面对国民党军的“围剿”，身处战场的中共中央不可能不把打破当前“围剿”当作头等重要的任务，这也是中共驻共产国际代表团在民族生存危机前思考方向逐渐发生变化而身处战地的中共中央却反应迟钝的重要原因。


在对形势作出严重判断的基础上，如何应付国民党新一轮的“围剿”，成为摆在中共面前的一道难题。


数年的苏区发展过程中，中共已经创造出一整套的运动战、游击战的战略战术，而且在每一个成功发展的苏区，我们都可以看到这些战略战术的总结，中央苏区军事与政治的结合尤其成功。如当时有人所总结的：

\begin{quote}

在游击斗争中，在群众工作上，必须采取毛泽东深入群众的斗争；但另一方面我们还须在军事上采取彭德怀的方式，游击队神出鬼没的袭击敌人，骚动敌人，分散敌人的目标，使其不知我们的踪迹，不易对付我们，在敌人注意我们游击队行动中，更利于群众斗争的发展与深入。\footnote{《中共厦门中心市委给漳属游击队信（1932年3月）》，《福建革命历史文件汇集》甲9册，第132页。}
\end{quote}


这一套在实践中得到提炼的军事、政治的成功经验，对中共展开反“围剿”战争有着重要参考价值。中共在第五次反“围剿”初、中期的战略指导中，对此也给予了高度重视，并坚持予以贯彻。1933年8月28日，中革军委下达密电，规定各省军区和军分区的区域和任务，要求各军区及所属部队配合主力红军作战，“进行机动的游击的战斗，以钳制麻醉疲惫迷惑敌方和瓦解当前之敌”。\footnote{《中央革命军事委员会密令（密字第一号）（1933年8月28日）》。}29日，周恩来在红一方面军政治干部会议上作《粉碎敌人五次“围剿”中中央红军的紧急任务》的报告，这是红军进行新一轮反“围剿”战争准备的主要文件。报告提出反“围剿”作战的主要方针：（一）猛烈地扩大和巩固红军，加强部队尤其是新战士的政治教育，发展党员，建立十人团以反逃亡。（二）加紧学习与提高军事技术，不仅要加紧练习打靶、投弹和白刃战等技术，而且要努力学习防空防毒、阵地战、堡垒战和使用新武器等技术，了解掌握一切战术原则和战斗动作，“不应忘记”游击战术。（三）有计划地建立各方面的正式战线与游击战线，开展游击队的游击战争，破坏国民党军的战争准备和供给。（四）各部队应在整个战略的意图下执行战斗任务，同时，“要能适应情况变动，在不违反整个意图下，发扬其机断专行的能力”，“要特别提倡活泼机动的运用战略战术，而不需要呆板的服从命令”。（五）加强对白区白军和红军中的政治工作，克服“敌人堡垒是无法打的”、“敌人封锁中的一切困难是无法克服的”等悲观失望情绪，同时反对胜利冲昏头脑的思想，“保证红军中思想的一致”。\footnote{周恩来：《粉碎敌人五次“围剿”中中央红军的紧急任务》，《斗争》第24期，1933年8月29日。}


这一时期，中共中央的一系列指示都贯串着运动战的思想。1933年8月，项英明确警示闽浙赣苏区领导人：“运动战是内战中主要方式，望努力学习。”\footnote{项英：《对闽浙赣工作的指示（1933年8月26日）》，《项英军事文选》，中共中央党校出版社，2003，第126页。}10月，中共中央致函闽浙赣省委，严厉批评单纯防御的思想，强调：“消极的堵防的政策，不论军事上与政治上都是有害的。在军事上，在部队数量暂时我们还比敌人小得多的时候，分兵把口，实际上就是使我们在敌人的残酷的集中兵力的进攻面前解除自己的武装。”指出：“最好的巩固苏区的办法，就是积极的开展深入白区的游击战争，发展与扩大苏区”。\footnote{《中共中央给闽浙赣省委信（1934年10月4日）》，《闽浙皖赣革命根据地》（上），第653～654页。}共产国际对运动战原则也反复强调。早在1933年3月，共产国际致电中共中央指示军事方针时明确谈道：

\begin{quote}

在保卫苏区时，对于中央苏区来说特别重要的是保持红军的能动性，不要以巨大损失的代价把红军束缚在领土上。应该事先制定好可以退却的路线，做好准备，在人烟罕至的地方建立有粮食保证的基地，红军可以在那里隐蔽和等待更好的时机。应避免与敌大量兵力发生不利遭遇，要采取诱敌深入、各个击破、涣散敌人军心和使敌人疲惫的战术，还要最大限度地运用游击斗争方法。\footnote{《共产国际执行委员会政治书记处给中共中央的电报（1933年3月19～22日）》，《共产国际、联共（布）与中国革命档案资料丛书》第13册，第353页。}
\end{quote}


4月，共产国际东方部负责人米夫发表文章阐述红军的作战方针：“假使要牺牲很多的战士时，红军绝不固守一个地方。红军避免与集中的大部敌人接触，引诱一部分敌人到苏区内来创造进行决战的胜利环境，而敌人则不得不在恶劣的环境下作战……红军还用游击战争与革命农民对政府军队士兵的影响来削弱和动摇敌人。红军在分散敌人的集中后，从侧面及后面攻击孤立的队伍，各个击破，消灭敌人。”\footnote{米夫：《中国革命危机的新阶段》，原载《共产国际》1933年4月号，苏区中央局《斗争》第23期转载，1933年8月22日。}1933年8月，中央苏区中央局机关刊物《斗争》刊发该文。9月和次年2月，共产国际又两电中共中央，强调运动战的作战方针：“中央苏区的主力不应参与阵地战，它们应该进行运动战，从两翼实行夹击。中央苏区要有预备力量，以对付任何突然袭击。”\footnote{《共产国际执行委员会政治书记处政治委员会给中共中央的电报（1933年9月29日）》，《共产国际、联共（布）与中国革命档案资料丛书》第13册，第509页。}“我军常在运动战中而不是突破敌人设防地区的战斗中取得很大胜利。应该充分利用我军的主要优势，即它机动作战和从翼侧突然进攻敌人的能力，而游击队同时从后方进攻敌人。”\footnote{《共产国际执行委员会政治书记处政治委员会给埃韦特、施特恩和中共中央的电报（1934年2月11日）》，《共产国际、联共（布）与中国革命档案资料丛书》第14册，第79页。}共产国际的这些意见，起码在理论上统一了中共的基本作战理念，李德回忆：“至于阵地战，不管是什么形式都不合适，这一点我们大家是完全清楚的。”\footnote{〔德〕奥托·布劳恩：《中国纪事1932～1939》，第63页。}苏区军政领导人普遍认为：“死守堡垒，对于我们是致命的打击”；\footnote{张闻天：《闽赣党目前的中心任务》，《斗争》第71期，1934年9月7日。}“要反对滥做支点单纯防御的堡垒主义”。\footnote{陈毅：《几个支点守备队的教训》，《革命与战争》第4期，1934年5月18日。}中共的这种作战方式，其对手蒋介石体会最深，他在1933年12月谈道：“现在我们打土匪，真面目的阵地战很少，而随时遭遇的游击战特别多。”\footnote{蒋介石：《为闽变对讨逆军训话——说明讨逆剿匪致胜的要诀》，《先总统蒋公思想言论总集》第11卷，第631页。}


游击战在各级、各地的有关指示中多次被提及。早在1933年8月，中共中央在致红四方面军信中指出“极大的注意开展游击战争，创造新的游击区域，是主力红军在决定胜负战争中取得胜利的必要条件之一。”\footnote{《中央致红四方面军的信（1933年8月25日）》，《中共中央文件选集》第9册，第315页。}中共领导人博古在关于第四次反“围剿”经验教训的总结中也不忘强调：

\begin{quote}

在粉碎敌人四次“围剿”中，游击战争的开展，非常不够，没有充分组织游击部队，有目的的配合红军作战，尤其没有在敌人的后方侧面发动新的游击区域，使敌人在四面受敌的形势中。譬如说，在北方战线上东黄陂战役的时候，我们能够在丰城东乡一带有一支像从前朱毛彭黄的游击队伍，那我们的胜利，不定更要比现在大得多。\footnote{博古：《为粉碎敌人的五次“围剿”与争取独立自由的苏维埃中国而斗争》，《秦邦宪（博古）文集》，第211页。}
\end{quote}


11月，中革军委代主席项英要求：“应在敌人后方要道，发展井冈山时代游击袭击的精神、第四次战役的挺进成绩，来配合作战以及转变战局。”\footnote{项英：《挺进游击队的任务（1933年11月6日）》，《项英军事文选》，第251页。}同月，中共粤赣省党代表大会政治决议强调在粤赣省普遍开展游击战争的重要性，同时批评各地不重视游击战争的倾向：“分兵把口还是普遍的存在着（在寻安信）。因此，更不能抓紧许多便利的条件，灵活的、广泛的在敌人侧后发展游击战争，创造游击区域。”\footnote{《中共粤赣省委政治决议案》，《闽粤赣革命历史文件汇集（1933～1936）》，中央档案馆编印，1984，第26页。}1934年1月，红军总政治部就开展游击队工作发出训令，指出“游击战争的开展，是主力红军在决定胜负的战争中取得胜利的必要条件之一”，要求“对任何忽视游击队工作的倾向，坚决进行斗争”。\footnote{《1934年1月5日总政治部训令，关于游击队工作》，《斗争》第42期，1934年1月20日。}4月，中共中央和苏维埃中央政府指示各地、各级党和苏维埃，“我们要以更多的地方部队，发展广大的游击战争”。\footnote{《中共中央委员会、中央人民委员会给战地党和苏维埃的指示信（1934年4月24日）》，《斗争》第58期，1934年5月5日。}以苏区中央局名义发表的社论也指出：“在敌人的堡垒政策面前，发展游击战争，可以使敌人力量很大的分散与削弱，使主力红军的战斗得到更便利的条件。”\footnote{《五次战役第二步的决战关头和我们的任务》，《斗争》第58期，1934年5月5日。}


值得注意的是，对当时的中共中央有着重要影响的王明在这一问题上认识也是清楚的。1933年12月，他在共产国际执委会第13次全会发言谈道，中共中央正在执行的军事政策的一个重要部分是：“反对那种‘左’的冒险的倾向，其具体表现，就是完全否认有暂时地部分地军事策略上的退却底可能和必要（例如，为避免和敌人过大力量底作战；或者为的抽出时间准备和找到好的机会再作进攻等等），就是对于保存红军实力有第一等意义这一点估计得不足或根本不了解，机械地了解巩固根据地任务，甚至以为可以牺牲大量红军力量去达到这一目的。”\footnote{王明：《革命、战争和武装干涉与中国共产党底任务》，《王明言论选辑》，第362页。}本着这一思路，他强调，打破蒋介石新一轮“围剿”的主要办法应是“把中央苏区的红军和游击队从内防守的军事行动，与红军及游击队在蒋军后方和两翼方面实行广大的游击战和运动战配合起来”；\footnote{王明：《新条件与新策略》，莫斯科，1934，出版者不详，第21页。}指出：“中国革命是一种长期性的坚（艰）苦斗争，六次‘围剿’，虽然是革命与反革命之间的残酷斗争的严重的步骤，然而它不仅不是最后决定中国命运的斗争，并且也不是决定胜负的斗争。”\footnote{《王明、康生致中共中央政治局函（1934年8月3日）》，转见周国全等《王明评传》，安徽人民出版社，1989，第255～256页。}


不过，共产国际、王明虽然注意到运动战、游击战的作战原则仍有其不可移易的效应，但远在数千里之外的他们其实并不一定对中共新的遭遇有深切的了解。应该看到，战争毕竟是敌对双方战略战术的相互较量，第五次“围剿”中，国民党军出动兵力较第四次“围剿”几乎翻了一番，直接参战兵力达到50万人左右，与此同时，红军总兵力增加不多，仅10万余人，\footnote{《赣粤闽湘鄂红匪实力调查表》，《赣粤闽湘鄂北路剿匪军第三路军五次进剿战史》（上），第三章，第7页。}在兵力上“我与敌常为一与四与五之比”\footnote{《中央对今后作战计划的指示（1933年6月）》，《中共中央文件选集》第9册，第226页。}。人数、武器装备和补给的绝对优势，使国民党军事实上可以主控战场。因此，当他们采取稳扎稳打方针，将“极力剥夺红军进行一般的运动战，尤其是遭遇战斗、袭击的机会”作为其战略战术“最基本的要点之一”\footnote{林彪：《短促突击论》，《革命与战争》第6期，1934年7月。}时，中共以往最为得意的运动战、游击战都面临严峻考验。以诱敌深入为例，由于国民党军力量上的绝对优势和进兵时的极端谨慎，导致红军集中优势兵力对敌方单兵突进部队实施打击的机会大大减少，李德谈道：“至于在我们区域内进行歼灭战的有利条件，只要不能诱敌深入，也就是说敌人不放弃堡垒战，那就没有希望得到……我们埋伏在这里，而敌人就可以丝毫不受干扰地继续推行它计划中的堡垒政策。这岂不意味着，我们自己放弃了苏区的重要地区，而不去利用时机歼灭敌人的有生力量。”\footnote{〔德〕奥托·布劳恩：《中国纪事1932～1939》，第60页。}而且，林彪还进一步指出：“诱敌深入的方法，在对付历经惨败而有无数血的教训的敌人，已经不是可靠的有效的方法了。”\footnote{林彪：《短促突击论》，《革命与战争》第6期，1934年7月。}对李德、博古军事指挥作了系统批评的遵义会议有关文件虽然强调第五次反“围剿”时，“诱敌深入的机会依然是有的”；但也清醒看到：“否认五次战争中敌人战略上的堡垒主义的特点，是错误的”。\footnote{陈云：《遵义政治局扩大会议传达提纲》，《遵义会议文献》，人民出版社，1985，第40页。}其实，早在第五次反“围剿”开始前，共产国际远东局在电报中就已经谈道：“总的说来，我们不能采取诱敌深入到苏区然后将其消灭的战术，因为要做到这一点，苏区还是太小了。”\footnote{《共产国际执行委员会远东局给共产国际执行委员会政治书记处政治委员会的电报（1933年4月3日）》，《共产国际、联共（布）与中国革命档案资料丛书》第13册，第374页。}所谓苏区太小了，当然不是说他们会无视相比前几次反“围剿”，苏区正在壮大的事实，而是基于其对外部环境变化的判断，即共产国际代表所说的：

\begin{quote}

我们的地区太小，不能遵循将敌人较大部队拖入陷阱的战术。我们应尽可能地进行大的战斗，或者在我们地区外围，或者在敌占区。由于敌人切割了我们的地区，控制着交通线和重要城市，我们不能让他们有机会将防御工事和基地向纵深推进。况且，如果敌人暂时占领我们的地区，他们在离开时会把它洗劫一空。在物质方面，这会削弱我们作战的能力，而在道义方面，我们离开团结一致的苏区，使之完全失去防御能力，就会损害农民对我们的信任。\footnote{《埃韦特给共产国际执行委员会的报告（1933年4月8日）》，《共产国际、联共（布）与中国革命档案资料丛书》第13册，第394页。}
\end{quote}


另外，由于国民党军推行坚固、密集的堡垒政策，红军运动战的区域、效用也大打折扣。中央红军反“围剿”初期在硝石、浒湾一带主动出击，试图在运动中把握机会消灭对方，由于对方处处设防、时时筑碉而难觅胜机。指挥作战的彭德怀、滕代远明确认为：“敌人正在大举集中的时候，利用堡垒掩护，使我求得运动中各个击破敌人机会减少，我军需要充实主力，储集力量与敌人大规模作战。”\footnote{《彭滕关于我军首先解决信河流域赵部的建议》，《中央苏区第五次反“围剿”》（上），第93页。}同样，湘赣红军在初期作战中遵照外线作战的原则向萍乡、宁冈等地的出击，也因该地“碉堡纵横，是湘军重兵驻屯的纵深地区，战场狭小，运动战机会很少”\footnote{萧克：《回忆湘赣红军》，《湘赣革命根据地》（下），第993页。}而遭失败。红十军在赣东北准备通过运动战集中力量打破国民党军进攻时，由于对方采取稳健推进战法，红军很难觅到战机，运动战往往被迫打成“攻击占领有利阵地做有工事的敌人”，\footnote{《刘畴西聂洪钧致朱德电（1934年4月20日）》，《闽浙皖赣革命根据地》（上），第710页。}运动战效能无法发挥。对于堡垒战中的双方对垒，彭德怀有一个形象的比喻：“等如猫儿守着玻璃里的鱼可望而不可得。”\footnote{《彭滕关于七日战况和改变行动的报告（1933年11月7日）》，《中央苏区第五次反“围剿”》（上），第99页。}可见，当时红军坚持其一贯的运动战原则确实遭遇到从未有过的困难。远在李德到达中央苏区之前，共产国际驻华代表及其军事顾问已经对困难予以充分注意，他们对困难的高度估计甚至导致其对国际方面的指示也并不认同，直指其会“助长错误的倾向”。\footnote{《埃韦特给共产国际执行委员会的报告（1933年4月8日）》，《共产国际、联共（布）与中国革命档案资料丛书》第13册，第394页。}


应该说，共产国际代表的认识并非无所根据，面对国民党军的战略改变，红军当然不应一成不变。如前所述，第三次反“围剿”战争胜利后，毛泽东就提出红军不应再采取防御式的内线作战战略，而应采取进攻的外线作战战略，也就是说，在新的形势下，红军不应再固守原来诱敌深入、内线作战的一贯战略，而应变内线作战为外线作战，主动出击，打到敌人后方去，从根本上破坏对手的部署。这和其原先的诱敌深入战术恰成对比，可谓兵无常形的最好注解。在《中国革命战争的战略问题》中毛泽东更明确总结了红军在不同阶段战略战术可能的变化：

\begin{quote}

红军的游击性，没有固定作战线，根据地的流动性，根据地建设工作的流动性，十年战争中一点也没有变化吗？有变化的。从井冈山到江西第一次反“围剿”前为第一个阶段，这个阶段中游击性和流动性是很大的，红军还在幼年时代，根据地还是游击区。从第一次反“围剿”到第三次反“围剿”为第二个阶段，这个阶段中游击性和流动性就缩小了许多，方面军已经建立，包含几百万人口的根据地已经存在。从第三次反“围剿”后至第五次反“围剿”为第三个阶段，游击性流动性更缩小了。中央政府与革命军事委员会已经建立。长征是第四个阶段。由于错误地否认小游击和小流动，就来了一个大游击和大流动。\footnote{毛泽东：《中国革命战争的战略问题》，《毛泽东选集》第1卷，第231页。}
\end{quote}


所以，毛泽东后来曾提出，福建事变爆发后，“红军主力无疑地应该突进到以浙江为中心的苏浙皖赣地区去，纵横驰骋于杭州、苏州、南京、芜湖、南昌、福州之间，将战略防御转变为战略进攻，威胁敌之根本重地，向广大无堡垒地带寻求作战。用这种方法，就能迫使进攻江西南部福建西部地区之敌回援其根本重地，粉碎其向江西根据地的进攻，并援助福建人民政府，——这种方法是必能确定地援助它的。此计不用，第五次‘围剿’就不能打破，福建人民政府也只好倒台”。\footnote{毛泽东：《中国革命战争的战略问题》，《毛泽东选集》第1卷，第218～219页。}而对反“围剿”中期由共产国际代表及军事顾问提出的进军湖南计划\footnote{参见《共产国际执行委员会政治书记处政治委员会给中共中央的电报（1934年1月2日）》，《共产国际、联共（布）与中国革命档案资料丛书》第14册，第7页。}也予以肯定，强调：“可以向另一方向改取战略进攻，即以主力向湖南前进，不是经湖南向贵州，而是向湖南中部前进，调动江西敌人至湖南而消灭之。”\footnote{毛泽东：《中国革命战争的战略问题》，《毛泽东选集》第1卷，第219页。}


放弃经营数年的中央苏区，集中兵力直捣敌之后方，这样的作战计划确实需要极高的想象力和极大的决心，它充分体现着毛泽东不羁的思路、傲岸的性格、特异的谋略，和其一贯的军事方针是相吻合的。但是，对于包括共产国际在内的军事、政治决策者而言，在对前景还没有完全绝望，又有前几次反“围剿”胜利的美好回忆背景下，确实不太可能一开始就直面这样一个破釜沉舟的设想。当时年轻的中共领导人很难承受由此带来的巨大风险，作出这种非常规的抉择，任何的闪失都有可能意味着需要承担严峻的政治后果，其中的压力，不是已经看到结果而又无须承担责任的我辈所可想象的。何况，以红军的现有力量，挺进到国民党政权的纵深区域，在一个不具有群众基础、回旋余地也有限的地区作战，风险也相当巨大。毕竟，毛泽东的根据地思想也正是建立在对方薄弱环节这一基础之上的。因此，中共中央当时选择在中央苏区展开反“围剿”作战确也不应感到意外。关键在于，在毛泽东看来，留在中央苏区继续作战，事实上也就意味着第五次“围剿”的不能被打破，而离开苏区进行外线的运动战，可以说是置之死地而后生的无奈选择。只是在第五次反“围剿”刚刚开始时，虽然人们普遍已经意识到了即将到来的严峻考验，但大部分人仍不会轻易认同这一估计。\footnote{事实上，毛泽东当时也曾表示：“苏维埃在全国民众对它的信仰日益增加中，在国民党及一切反革命派别的欺骗日益破产中，将坚决地粉碎六次‘围剿’，以便努力阻止帝国主义殖民地化中国的道路，努力争取在全国范围内苏维埃革命的胜利。”（毛泽东：《中华苏维埃共和国中央执行委员会与人民委员会对第二次全国苏维埃代表大会的报告》，《苏维埃中国》，中国现代史资料编辑委员会编印，1957，第252页）固然，不能把毛泽东这一讲话和其内心真实想法完全等同，这其中很可能有服从纪律的成分，但要断定他当时就有明确的失败认识，也未必符合事实，更多的可能是其积数年反“围剿”战争经验形成的敏感和嗅觉。}毕竟，“战争中一切行动追求的都只是可能的结果，而不是肯定的结果”。\footnote{克劳塞维茨：《战争论》第1卷，中国人民解放军军事科学院译，中国人民解放军总参谋部出版局，1964，第215页。}


当毛泽东肯定外线作战的思想时，他实际是把运动战的思路放到全国这样一个大的棋盘上作出重新的定位，这可以说是深得运动战之精髓，而共产国际和中共中央虽坚持运动战的原则，但限于苏区内部的运动战由于毛泽东所说的“游击性流动性更缩小了”，事实上难以发挥其曾经有过的威力。不过，正如遵义会议决议谈到的：“所谓运动战粉碎堡垒主义，即是在堡垒线内待敌人前进时大量消灭敌人的部队，在堡垒线外即是红军转到广大无堡垒地带活动时，迫使敌人不得不离开堡垒来和我们作运动战。”\footnote{《中央关于反对敌人五次“围剿”的总结的决议（遵义会议）》，《中共中央文件选集》第10册，第460页。}在内线利用敌人前进实施突击以消灭敌人毕竟是当时条件下运动战的两种形式之一，这一形式在当时的代表就是“短促突击”战术。

\section{“短促突击”战术}


中共在第五次反“围剿”作战中选择“短促突击”，经历了一个变化过程。反“围剿”开始前，远东局提出要使用常规部队和游击队打击对手后方，认为：“这种战术可以使我们在苏区之外顺利作战，避免重大损失。”\footnote{《共产国际执行委员会远东局给共产国际执行委员会政治书记处政治委员会的电报（1933年4月3日）》，《共产国际、联共（布）与中国革命档案资料丛书》第13册，第374页。}中革军委期望：“假若我们在很短的时间以内，能够在黎川获很大的胜利，那末，蒋介石的第二步骤就更要小心了，那么我们便可采取旋回政策，推迟和避免与他决战。”\footnote{《中央对当前作战方针的意见（1933年10月5日）》，《中央苏区第五次反“围剿”》（上），第78页。}这一点，其实在红军将领中也有相当共识，林彪、聂荣臻在战役开始前夕建议：

\begin{quote}

敌人采用逐步延伸的办法，首先完成在永（丰）吉（水）间之封锁线，然后再移其兵力，进行永（丰）乐（安）封锁线……在此情况下，如我军仍在现地不动，则有使敌封锁线完成可能，故我们意见，应以我们的行动调动敌人，以寻求击敌机会，以破坏敌封锁计划，我们除以一部正面钳制外，余应到北线敌人之后方翼侧及间隙中活动。\footnote{《聂荣臻年谱》（上），人民出版社，1999，第105页。}
\end{quote}


在这一思路指导下，反“围剿”作战初期，当时由项英代主席的中革军委曾尝试过到苏区外围进行外线作战，顶出去打，力争主动的战法，其攻击方向主要为国民党军较为薄弱的黎川、硝石地区，试图在此打破国民党军封锁线，打通与闽浙赣苏区的联系，破坏国民党军的堡垒战术。虽然硝石战役通过运动作战截断国民党军一部，使蒋介石不得不承认其“七十九师在硝石吃亏”，\footnote{蒋介石：《为闽变对讨逆军训话——说明讨逆剿匪致胜的要诀》，《先总统蒋公思想言论总集》第11卷，第629页。}但这一作战方针未能取得明显效果。由于国民党军兵力上的绝对优势和对防御工事的特别重视，红军进入国民党军区域后，处处受制，不仅难以调动对方，自身后路还常有被切断之虞。时任红十三师政治部主任莫文骅谈到：“从表面上看，我们是在主动进攻，试图在运动战中消灭敌人，实际上由于红军兵力不集中，又是在敌人堡垒密布的白区作战，故我们一开始就处于被动的境地，可说是虎落平川。”\footnote{《莫文骅回忆录》，解放军出版社，1996，第229页。}由于红军武器落后，攻坚能力不足，堡垒的威力不能低估。1934年年中，湘鄂赣红军试图从国民党军碉堡群中突围时，“把机枪集中向碉堡攻，打不烂，冲也冲不过去，牺牲很大。从早晨开始突围，一直打到下午快天黑，一道封锁线还没有突过”。\footnote{傅秋涛：《关于湘鄂赣边区内战中期后期历史情形报告》，《湘鄂赣边区史料》，中共中央宣传部党史资料室1954年印行，第77页。}


在难以顶出去打，又希望坚守苏维埃现有区域的情况下，红军不得不退而求其次，开始在苏区外围与国民党军展开运动防御战。当时，人们普遍认同：“消极的防御一定是失败的”，\footnote{华夫（李德）：《革命战争的迫切问题》，《革命与战争》第2期，1934年4月1日。应该指出，虽然红军的基本战斗方针是运动战，但一定条件下的防御战也并不违背红军的基本作战原则。毛泽东指出：“基本的是运动战，并不是拒绝必要的和可能的阵地战。战略防御时，我们钳制方面某些支点的固守，战略进攻时遇着孤立无援之敌，都是应该承认用阵地战去对付的。采取这样的阵地战制胜敌人的经验，我们过去已经不少；很多的城市、堡垒、寨子，被我们打开，某种程度的敌人野战阵地被我们突破。以后还要增加这一方面的努力，补足我们这一方面的弱点。我们完全应该提倡那种在情况需要而且许可下的阵地攻击和阵地防御。我们所反对的，仅仅是在今天采取一般的阵地战，或者把阵地战和运动战平等看待，这些才是不能许可的。”见毛泽东《中国革命的战略问题》，《毛泽东选集》第1卷，第230～231页。}“应采取积极的和运动的防御”。\footnote{项英：《关于战术问题的指示（1933年10月15日）》，《项英军事文选》，第222页。}对防御中应坚持运动战原则、力争主动，在理论上有清醒认识。被认为推行了消极防御政策的李德和项英当时都明确指出，红军反“围剿”战争中“最基本的要求，是要在敌人堡垒主义的条件下，寻求运动战，不要在进攻堡垒中，来消耗我们的兵力和兵器，要在堡垒外，于敌人的运动中，来消灭其有生兵力”；\footnote{华夫（李德）：《论红军在堡垒主义下的战术》，《革命与战争》第3期，1934年4月20日。}“应该避免进攻要塞堡垒地域，甚至避免正面进攻停止的敌人。我们战术的特质就是要搜求运动中的敌人，特别是他的翼侧施行迂回，或因地形和时间的关系施行包围，以及迅速而猛烈地突击敌人纵队第二、第三梯队的翼侧。”\footnote{项英：《关于战术问题的指示（1933年10月15日）》，《项英军事文选》，第222页。}国民党军在广昌战役后的总结中也谈道：“匪每欲以碉楼线，袭留我主力于正面，利用其重兵，袭击我之侧背。”\footnote{《赣粤闽湘鄂北路剿匪军第三路军五次进剿战史》（上），第八章，第54页。}国民党方面谈到的这种战法，就是中共在第五次反“围剿”中运用最多的“短促突击”战术。


所谓“短促突击”，即以一部防御吸引敌军，同时将主力埋伏在附近，当敌军出现在我前沿阵地时，以埋伏之主力部队“施行短促的突击及袭击，以便于堡垒前瓦解敌人”。\footnote{华夫（李德）：《革命战争的迫切问题》，《革命与战争》第2期，1934年4月1日。}李德具体规定了这一战术的几个主要原则：“向敌人运动中的部队进行短促的侧击”；“在敌人后续梯队或堡垒内来的增援队未到达前，迅速解决战斗”；“要最坚决的使用最高度的主力作战，以便确实的避免延长战斗”；“迅速转变自己的突击方向，主要的是利用敌人诸纵队的内翼侧，在其诸纵队间执行机动”。\footnote{华夫（李德）：《论红军在堡垒主义下的战术》，《革命与战争》第3期，1934年4月20日。}简言之，就是要吸引敌军于堡垒之外，集中优势兵力，迅速对敌军实施包围、速决歼灭，“他不是一个正规的攻击战，也不是一个正规防御的战斗，他是混用着各种复杂的战斗方法”，\footnote{林彪：《短促突击论》，《革命与战争》第6期，1934年7月。}包含了这一时期中共用兵的几个基本原则：迅速、突然、机动、集中兵力。


从目前材料看，这一战术发端于第五次反“围剿”前夕。1933年9月中旬，共产国际远东局给中共中央的电报中已经显露了这一战术的雏形：“我们应该让敌第一梯队部队向前推进几里。在这时我们将迅速在两翼运动，以便出人意料地迎击敌主力部队。在先歼灭敌第二梯队之后，再以小股部队击溃其第一梯队。”\footnote{《共产国际执行委员会远东局给中央苏区的电报（1933年9月12日）》，《共产国际、联共（布）与中国革命档案资料丛书》第13册，第493页。}第五次反“围剿”开始后不久，这一战术在实践中得到运用。1933年10月中旬，项英就硝石战役发出指示：“三军团应力图在十三、十四两日，向西及西南以个别的短促的打击在一师以内之先头部队，不应与敌之大兵力作战，不应向硝石作任何攻击。”\footnote{项英：《关于红三红五军团作战行动及任务的指示（1933年10月13日）》，《项英军事文选》，第217页。}11月底，项英再次提到：“对敌各个部队不大于一师（的），给以短促迅速的突击……要避免与敌人的兵力过早开始决战。”\footnote{项英：《对各部任务及动作的指示（1933年11月28日）》，《项英军事文选》，第290页。}周恩来、朱德则在致林彪、聂荣臻电中，就截击国民党军吴奇伟师指示：“这一截击应是迅速突然短促的，绝对不应正面强攻。”\footnote{《周恩来年谱1898～1949》，人民出版社、中央文献出版社，1989，第255页。}12月22日，周恩来、朱德又致电红军指挥员，强调红军应切忌正面强攻与相持恋战，要力求在敌移动中从侧翼进行最短促干脆的突击。\footnote{周恩来、朱德1933年12月22日致红三、五、七、九军团电，《周恩来年谱1898～1949》，第256页。}11月27日，“短促突击”已出现在红军指挥员的日记中，红五军团十三师师长陈伯钧明确写道：“我军有集结主力，以‘短促突击’侧击该敌之任务”。\footnote{《陈伯钧日记（1933～1937）》，上海人民出版社，1987，第128页。}“短促突击”战术在实战中运用已比较经常。


应该说，对“短促突击”战术作出最集中阐述的是代表共产国际驻中央苏区的军事顾问李德。李德总结了国民党军新一轮“围剿”中所取战略，指出国民党军在战略上放弃了过去的坚决的突击，而采取持久消耗的作战方针，面对这一新形势，红军像以前那样采取诱敌深入的大规模运动作战已不太可能，短促突击的方法应是相对可行的选择。短促突击主要包含着两方面的内容，一是利用支撑点的防御以吸引敌人。支撑点的防御不是目的：

\begin{quote}

运动防御的目的不是为要顽强扼守阵地或消灭敌人，而是为争取时间及诱引敌人（诱引敌人远离其堡垒以便我突击队隐蔽的突击之）；因此运动防御只应同敌人的先头部队（侦察、前卫游击队等）作战，并迫使敌人的主力展开，当执行了这个任务时，即有计划的转移到后一个地区……我们要记着，运动防御是为着保证我们主力在有利条件下施行突击的机动，如过早的退出战斗或顽强的战斗，都不能保证这些条件的构成。\footnote{华夫（李德）：《反对曲解我们的战术》，《革命与战争》第4期，1934年5月18日。}

发展游击战争，支撑地域的防御及进攻敌人的堡垒，为革命战争的辅助方式。这些战争方式，应协助造成战术的环境，使我们能实现基本的原则：即是以主动的机动，于堡垒外，消灭敌人的有生兵力。\footnote{华夫（李德）：《论红军在堡垒主义下的战术》，《革命与战争》第3期，1934年4月20日。}
\end{quote}


二是对被诱引的敌人在运动中实行突然的集中打击，歼灭敌人。这是短促突击的关键所在：“主力的机动和突击是有决定意义的”，“作战时应使用全力以便一举而迅速地解决战斗”。\footnote{华夫（李德）：《革命战争的迫切问题》，《革命与战争》第2期，1934年4月1日。}


李德的解释，和林彪专指对脱离碉堡的敌军的突击更为宽泛，实际上，红军一贯的诱敌深入的运动战战术其关键词也不能离开短促、突击，从这一点上说，短促突击战术在战术原则上并不具有多少新的内容，只是在国民党军新的作战原则下运动战的策略调整。\footnote{关于运动战，毛泽东指出：“作为战争内容的战略内线、战略持久、战略防御中的战役和战斗上的外线的速决的进攻战，在战争形式上就表现为运动战。”（毛泽东：《论持久战》，《毛泽东选集》第2卷，第497页）中央苏区时期根据苏联军事著作给予运动战的定义是：“‘运动战’，就是对于阵地战而言的话，通常就是指发生于兵团运动中而不是长久固着一地的战斗，例如遭遇战斗，追击，退却，以及在通常之防御地带的进攻及防御等。然而，在大兵团的运动中，其一部分有作阵地战的。”（冰澈：《军语解释》，《革命与战争》第1期，1932年8月1日）运动战的特点是“正规兵团，战役和战斗的优势兵力，进攻性和流动性”，同时也包括“辅助作用的阵地战”、“‘运动性的防御’和退却”。（毛泽东：《论持久战》，《毛泽东选集》第2卷，第498页）以此对照“短促打击”，应该说，两者基本原则具有相当的一致性。}


总体来看，短促突击战术，作为持久防御总方针下的一个战术原则，着眼于防御战中尽可能发挥红军善于集中力量打运动战的优势，以歼灭敌军有生力量而不以单纯保有地域为原则。这一战术要求尽量机动掌握兵力，避免与敌人过多地消耗，还要尽可能地减少自己的牺牲。中革军委强调：“在我们的条件之下，战斗的胜利不是占领地方，而是消灭敌人的有生力量及夺取其器材。”\footnote{《中革军委关于十月中战役问题致师以上首长及司令部的一封信（1933年11月20日）》，《中央苏区第五次反“围剿”》（上），第112页。}“要以最高度的保持有生力量，物质基础及我们新的原则为出发点。”\footnote{项英：《对各部任务及动作的指示（1933年11月28日）》，《项英军事文选》，第290页。}李德也谈道：

\begin{quote}

如遭遇的突击未成功，而敌人又已构成“战斗正面”时，则不宜继续强攻固守的敌人，而应当退出战斗，争取其他方向的先机之利。立于主动地位，决心和实现决心的灵活性，在这里是有重大意义的。由已得的胜利争取全部胜利的顽强性与在不利时勇敢退出战斗，并不是矛盾的，而是相互辅助的，这就在乎良好（沉着与坚决）指挥员的适当运用。\footnote{华夫（李德）：《再论战术原则》，《革命与战争》第4期，1934年5月18日。}
\end{quote}


在这一精神指导下，最大限度地杀伤对方而尽可能保持己方有生力量仍然被视为红军极其重要的战斗原则，红十三师在德胜关一带阻击国民党军遭受较大损失后，遭到军团首长的严厉批评，责其：“对贯彻军委关于极力保障物质基础和有生力量的指示万分不够，甚至是罪恶。”\footnote{《陈伯钧日记（1933～1937）》，第142页。}


除李德外，红军主要指挥员朱德、彭德怀、林彪等也曾撰文论述短促突击战术。彭德怀谈道：“在敌跃进和推进时，灵活的运用攻击的战术动作——侧击和短促的突击，来取得敌人资材，根据自己的特长和敌情，是有可能的；而且也只有这样积极动作，争取每次战斗的胜利，才能展开战役上的胜利，完成持久战略。”\footnote{彭德怀：《顽强防御与短促突击——给某师长的信》，《革命与战争》第9期，1934年9月10日。}朱德在总结高虎脑战斗时指出：“突击队能英勇与适时施行短促突击，守备队能坚决与顽强地抵抗互相配合，是有很大意义的。”\footnote{朱德：《一个支撑点和短促突击的战例》，《革命与战争》第9期，1934年9月10日。}林彪、聂荣臻1934年2月致朱德要求实施运动战的建议中也谈道：

\begin{quote}

我主力所在地域如附近有敌，则诱敌和放敌大踏步，以便我主力在敌运动中消灭之，如我主力不在某地而该地有敌前进时，则应以一部兵力进行运动防御战，滞敌前进……我主力军到达后，如见敌工事尚不坚固，则以主力攻击之，如敌堡垒已极坚固，但联络堡尚未做好，则应佯攻与围攻其堡垒，而打击敌来作联络堡之部队或增援队。\footnote{《用运动战消灭敌人》，《聂荣臻军事文选》，解放军出版社，1992，第28页。}
\end{quote}


这一建议实际包含了林彪后来发表的阐述“短促突击”文章的基本思想。


在总体战略受制于人的背景下，短促突击虽然可以发挥一定的效果，但成绩仍然有限。首先，在国民党军采取稳扎稳打战略的背景下，要抓住对方突进部队实施突击机会十分有限，如李德自己所说的，国民党军大胆前进，“一下子离开其基本堡垒十里十五里至二十里”的情况，往往发生在“二○至三○里的地域上集中十个师以上，而在十里以内的地域内有三个到四个师的突击队”\footnote{华夫（李德）：《再论战术原则》，《革命与战争》第4期，1934年5月18日。}的前提下，在对方兵力如此厚集时，红军要想取得歼灭战的战果，困难重重。其次，国民党军“推进的距离更加短促，力求避免其翼侧暴露缘着其预定的道路两侧推进，其正面很窄狭，以极大纵深集团推进”，\footnote{彭德怀：《顽强防御与短促突击——给某师长的信》，《革命与战争》第9期，1934年9月10日。}红军突击即使抓住其部分部队，也难以形成歼灭战。林彪、聂荣臻当时就谈道：“在敌人堡垒外的近距离或从堡垒间隙中去求运动战，结果仍变成堡垒战，以大部队在这种场合想行短促而突然的突击，结果打响之后仍然不易摆脱。”\footnote{《聂荣臻年谱》（上），第114页。应该指出的是，他们强调如上问题的目的并不在于否定“短促突击”，而是希望“应于敌前进而于敌之运动中和初到时消灭之……只要敌后尾离开其后面堡垒在十里以外，我们是有把握消灭他的”。}再次，红军在兵力、武器均处严重劣势情况下，短促突击在短兵相接这一作战阶段，损失和牺牲仍嫌过大，红军难以长期承受。因此，短促突击战术可以说是在战略被动的大背景下一种无奈的战术选择，它虽有其成立的背景和合理性，但不像中革军委和李德所期望的那样可以发挥出巨大的效能。正如林彪所指出的：

\begin{quote}

我们在战略战术上，是一方面要极力利用革命战争的各种辅助方式（游击战争防御等）；但最基本是要用巧妙的机动，以主力寻求在更宽大无堡垒的地域，进行正规的，大模范的运动战，大量的消灭敌人的有生力量。短促突击虽然也是运动战的一种，但我们如完全束缚在这种战斗方式中则是非常错误的。在另一方面，如果不相信这种短促突击的战斗有消灭敌人的可能，有造成战役上胜利的可能，而忽视这种战斗，则更是危险的损害的。\footnote{林彪：《短促突击论》，《革命与战争》第6期，1934年7月。}
\end{quote}


林彪这一段话，颇值重视，其以主力进到更宽大无堡垒地域进行运动战的想法，和毛泽东的一贯思路是一致的，但由于现实条件的限制，这一思路暂时难以付诸实施，现实的可能仍是在堡垒线内寻找运动作战的机会，这是他支持“短促突击”的基本理由，也是该战术成立的不容忽视的背景。


其实，不应否认，当时中共军事指挥者一直也在探索更多的应对方略。他们要求各级干部：“在任何情况下，不应呆板机械的执行指示和命令，而应深刻了解其意旨，并依所受命令的意旨及实际的情况，勇敢机断地专行起来。因此，必须经常地估计敌情、我军、地形和时间，当每一情况变移中，应即速定下自己基本的决心。”\footnote{朱德、周恩来、王稼蔷（祥）：《关于支撑点的构筑和防御以及对付敌人堡垒之战术的补充指示（1934年2月25日）》，《赣粤闽湘鄂北路剿匪军第三路军五次进剿战史》（下），附件（五），第20页。}李德进一步指出：“最危险的就是简单化的及机械的应用战术原则。敌人和我们的战术都是在发展中变更中成就中，若果以这些原则引以为足时，那就要在目前的战斗环境中算落伍了。”\footnote{华夫（李德）：《反对曲解我们的战术》，《革命与战争》第4期，1934年5月18日。}他强调：“革命军队的基本优点，高度的机动性独断专行以及勇敢的突击。”\footnote{华夫（李德）：《革命战争的迫切问题》，《革命与战争》第2期，1934年4月1日。}广昌战役开始前夕，林彪写信给中革军委，指出：

\begin{quote}

有些重要的负责同志，因为他以为敌人五次“围剿”中所用的堡垒政策是完全步步为营的，我们已失去了求得运动战的机会，已失掉一个战役中消灭几个师的机会。因此遂主张我军主力分开去分路阻敌，去打堡垒战，去天天与敌人保持接触，与敌对峙……这种意见我是不同意的。

我觉得我们主力通常应隐蔽集结于机动地点，有计划的尽可能造成求得运动战的机会，抓紧运动战的机会，而于运动战中我主力军大量的消灭敌人，每次消灭他数个师。\footnote{林彪：《关于作战指挥和战略战术问题给军委的信（1934年4月3日）》，转见黄少群《中区风云——中央苏区第一至第五次反“围剿”战争史》，第326页。}
\end{quote}


这实际是主张采用更灵活、广泛的思路来应对新的战争。

\section{红军的正规化建设和防御原则}


在第五次反“围剿”准备过程中，红军的正规化建设受到较多的重视。应该指出，随着红军的逐渐壮大，苏维埃区域逐渐稳固，红军正规化建设得到更多强调确属顺理成章。博古在第五次反“围剿”前夕对红军将士提出要求：

\begin{quote}

第一，我们应该最好的来使用我们现有的武器，射击、刺杀、高射、手榴弹之掷抛、机关枪自动步枪之熟练的使用，火力之配备及协同动作，基本动作之纯熟，警戒的严密，行军、夜间动作、冲锋等等，必须使得我们在这些方面，我们应该发扬到最高限度，最好的最精确的熟练的使用我们一切现有的武器与技术。第二，我们必须最好的来防御一切敌人的新式武器与夺取他们来武装自己。譬如防空对空射击、防毒、坦克车、装甲浅水兵舰等等，在五次“围剿”中必定将更多更大的采用这种新式武器，我们必须使每一个指挥员战斗员都了解这些武器作用与效能以及抵抗他的方法，不致于在它突然出现时，给我们以慌乱或重大的损失，并且要夺取这些武器为我们使用。第三，对于指挥员应该是更大的加深战术与战略的研究，要最灵活的运用我们在国内战争中丰富的经验及最新战略与战术的原则。\footnote{博古：《提高我们的军事技术（1933年8月13日）》，《秦邦宪（博古）文集》，第223页。}
\end{quote}


在苏维埃运动中享有很高声誉的领袖方志敏也提出：“游击主义的残余，什么事都随便、马虎、不认真，以军事为儿戏，这是我们创造铁的红军的工作中最有害的残余。”\footnote{志敏：《建设我们铁的红军》，《突击》第7期，1934年2月。}其实，毛泽东对红军的游击主义和正规化建设曾有中肯评价，他认为：“游击主义有两方面。一方面是非正规性，就是不集中、不统一、纪律不严、工作方法简单化等。这些东西是红军幼年时代本身带来的，有些在当时还正是需要的。然而到了红军的高级阶段，必须逐渐地自觉地去掉它们，使红军更集中些，更统一些，更有纪律些，工作更周密些，就是说使之更带正规性。在作战指挥上，也应逐渐地自觉地减少那些在高级阶段所不必要的游击性。在这一方面拒绝前进，固执地停顿于旧阶段，是不许可的，是有害的，是不利于大规模作战的。”\footnote{毛泽东：《中国革命战争的战略问题》，《毛泽东选集》第1卷，第232页。}这都正面肯定了一定条件下红军正规化建设的必要性。


根据红军逐渐正规化的原则，第五次反“围剿”开始前，红军指挥部对红军编制作了若干整顿：“缩小部队编制，由七千余减为五千余人、枪由三千八百减为三千五百为一师。充实战斗员，大大减少非战斗员和杂务人员，以适合于山地战和干部能力。”\footnote{项英：《关于部队缩编及逃亡问题（1933年8月30日）》，《项英军事文选》，第145页。}营、连、排、班的编制也予以精简、充实，每连九个步兵班，一个轻机枪班，每班9人。这一整顿原则和国民党军一样，都是精简编制，便于部队运动，适合在山地作战。显然，红军加强正规化建设、提高正规作战的能力并不等于说否定运动战、游击战的方针，而是要求部队更统一、集中，纪律更严明，作战能力有更大的提高。苏区中央局在答复湘赣省委的报告时明确指出：“你们在报告上提到要消失（灭）游击战争的残余，或游击的残余，这是错误的。对游击战争我们不但不反对，并且在革命战争中时时要运用他。”\footnote{《苏区中央局给湘赣省委信（1932年3月4日）》，《湘赣革命根据地》（上），第233页。}反对的只是“游击主义的残余（不爱惜武器忽视军事技术等）”。\footnote{萧华：《以战争胜利来纪念八一》，《青年实话》“八一”增刊，1933年7月30日。}


第五次反“围剿”开始前后，中革军委之所以更多强调正规化，和国民党军的逼迫直接相关。1933年3月，中共中央在给鄂豫皖省委的信中强调：

\begin{quote}

在从前的时期，为建立并巩固苏区以反抗敌人，旧的游击战争是充分的，因为当时的敌人还没有很好的准备，用很大的力量向我们进攻。但只是这个策略，决不能破坏敌人强大的力量，决不能防止敌人在我们苏区周围筑起堡垒与交通网。起初，我们军队还缺乏强大而有组织的单位，便使我们在当时难于做正式的战争。但在这数年之间，我们的武装力量，不仅在数量上增加，而且军备和组织上也大大进步了，游击战与正规战的适当配合，已为可能之事；这个配合作战，更因为敌人之变更策略而成为必需。\footnote{《中共中央致鄂豫皖苏区党省委信》，《鄂豫皖革命根据地》第1册，第190页。}
\end{quote}


周恩来在1934年初也指出：“我们现在已开始进行堡垒战、阵地战、夜间战斗。这是实际战斗的要求使我们走上了这一步，自然我们主要的作战形式还是运动战。但目前已常常可以看到，遭遇战、运动战很快就转成阵地战……自然，我们要求从运动战中来消灭敌人，这是最有把握的。但是敌人不会总是那样蠢笨，我们要估计到各种情况中的战斗。”\footnote{周恩来：《一切政治工作为着前线上的胜利（1934年2月7日）》，《周恩来军事文选》第1卷，第322页。}


从国共两军当时的对垒状况看，面对国民党军日渐强盛的进攻，红军在被迫进行的对垒作战中，确实面临着加强部队整顿、训练及提高作战能力的任务，当时有关战报清楚地显示出红军作战能力、素养的欠缺：

\begin{quote}

在冲锋时及击退敌人而占领其阵地后，和敌反冲锋到来时均未发扬火力。指挥员忘记了运用机关枪，没有指示机枪的射击目标，及阵地机枪在战斗中随便摆在阵地上无人过问，枪口有的还朝着后方（二师）。全师机枪只两枝枪带了水，其余均未带水……部队对随时准备战斗的指示是不充分的，对于火力的运用是不注意的，战斗中火力与运动的配合是差到惊人的程度。

战斗中所表现一股作气的自发的勇敢是很好的，开始冲锋是很迅速的、勇敢的，但一遇到敌人较顽强的抵抗时或较有力的火力阻拦时，就在敌前停滞起来了，也不发扬火力也不跃进，大家挤在一堆让敌人最有效的火力射击和遭受敌人的手榴弹的掷炸，死伤枕藉。等到敌人的反冲锋到来，敌人走近我们的人堆附近以集束的炸弹投入我们的人丛间，一个炸弹就能炸着我们几个人。\footnote{《一军团云盖山大雄关沙岗上等地附近战斗经过详报（1933年11月15日至19日）》，《中央苏区第五次反“围剿”》（下），第36～37页。}
\end{quote}


相对于国民党军的整训，当战争爆发，两军直接对垒时，看得出来，红军在这方面做的工作不是多了，而是仍然不够。


不可否认，第五次反“围剿”期间，由于当时对垒双方的特殊态势，中革军委在坚持运动战总原则同时，对一定程度的防御给予了更多的重视。客观看，面对国民党军的步步进逼，红军从事一定程度的防御战应为势在必然。1933年上半年，周恩来就曾谈道：“照战术原则说来，在政策上或战术要点上，构筑永久或半永久筑城的要塞，而以极少数兵力编成的要塞军孤立死守，其目的不仅要使野战军行动自如，并且要使野战军利用其限制敌军行动的良机，遂行其企图，并且依此保护广大苏区，使敌人不敢突入，更以之为据点而向外发展，这就我军设备要塞说来是应如此。”\footnote{周恩来：《论敌人的要塞》，《红色战场丛刊》，工农红军学校1933年6月编印，第80页。}10月，中革军委命令：“各作战分区须依照军委密令在战略与战术的必要上立即选择险要地，并完成必要的土围和工事，准备相当粮食与弹药，具体指定相当地方部队据守。”\footnote{《中央革命军事委员会紧急命令（1933年10月5日）》，《红星》第10期，1933年10月8日。}不过，按当时军事领导人的想法，这种防御并不等于死守，李德指出：“防御时应布置积极的防御，以少数兵力及火器守备堡垒，而主力则用来施行短促的突击及袭击，以便于堡垒前瓦解敌人。”“在某一方面集中主力以行坚决的突击并在堡垒外消灭敌人的有生兵力。游击战争和防御虽是革命战争必须的方式，但只是辅助的方式，主力的机动和突击是有决定意义的。”\footnote{华夫（李德）：《革命战争的迫切问题》，《革命与战争》第2期，1934年4月1日。}周恩来也强调：“我们的防御是攻势防御，它不能将苏区周围都修起像万里长城的支撑点来守备，这并不是由于红军数量尚少于敌人，而是因为这种防御是为着配合进攻而防御。我们反对单纯防御路线，要进行运动的防御，要进行钳制敌人以便于突击和消灭敌人的防御。”\footnote{周恩来：《一切政治工作为着前线上的胜利（1934年2月7日）》，《周恩来军事文选》第1卷，第321页。}当然，在反“围剿”战争的实际进程中，由于红军很快陷于被动防御，周恩来等强调的攻势防御原则并不总能真正执行，消极防御的战斗不乏其例。其中原因十分复杂，不能简单归咎于某一个体或团体，正如任弼时后来总结这段历史时谈到的：“我们也不能把所有的筑堡垒的事情都归之于新路线。”\footnote{任弼时：《在湘赣工作座谈会上的发言（1944年10月26日）》，《湘赣革命根据地》（下），第830页。}


随着红军第五次反“围剿”的屡战不利，从1934年年中开始，中共开始将撤离中央苏区，实施战略大转移列入议事日程，此后中共的作战方针某种程度上就是为这一全局性计划的实施争取时间了，被动防御成为这一时期中共作战的主流。期间，红军曾经策动过以“六路防御”对付“六路进攻”，将顽强防御作为阻止国民党军迅速深入苏区的唯一办法。即使考虑到红军战略转移争取时间的需要，这种战略指导仍然显得过于呆板，失却了此前红军机动灵活的特点，这应该和连续失败严重打击中共领导层及红军指挥部信心不无关系。同时，这一时期由于共产国际军事顾问李德个人军事专权严重，对发挥指挥员能动性注意不够，经常根据不准确的情报遥制前方战事，更带来了红军作战上的困难。彭德怀后来谈道，当时指挥部甚至“连迫击炮放在地图上某一曲线上都规定了。实际中国这一带的十万分之一图，就根本没有实测过，只是问测的，有时方向都不对”\footnote{彭德怀：《彭德怀自述》，第191页。}。广昌战役前，他已就此坦率提出意见：

\begin{quote}

在我们历次战役中，感觉着我们战略决心的迟疑。战略战术的机动，还未发展到应有的程度，失掉了许多先机，使应得的胜利，推迟下去，或使某一战役变成流产。而在战术上企图挽救过时的战略动作，结果把战术动作限制得过分严格，失掉下级的机动，变成机械的执行。（对）每一分钟的敌情变化和某地带地形的特点，不能灵活机断专行完成所给予的任务。这一方面，应拿许多适当例子来发扬指挥员的机动，加强其战术的修养；另一方面，每一次战斗，应给予其总的任务和各个的任务，不宜限制他执行的机动。\footnote{《彭德怀关于作战问题给军委的信（1934年4月1日）》，《中央苏区第五次反“围剿”》（上），第153页。}
\end{quote}


彭德怀的意见，可以从当时红军前线指挥员的亲身感受中得到印证。陈伯钧日记记载，在广昌南作战中，军委“判断敌人要南进，必须先攻占鸡公脑（这是照军委发下之地图），所以部署重新变更……这一部署主要是侧重鸡公脑，但实地的鸡公脑与地图上不合……给军团首长一报告，估计敌人不会绕攻鸡公脑，但军委因按错误的地形图指导作战，所以有此决心”。\footnote{《陈伯钧日记（1933～1937）》，第252页。}这些，确实在相当程度上影响到红军战斗力的发挥。


不过，正是彭德怀，在高虎脑战斗后曾提出过另外一种想法，给我们理解那次战争提供了一种独特的思路。他谈道：

\begin{quote}

敌人战术因其不愿意脱离堡垒之故，（一）不愿意行较大的翼侧机动而只以纵深梯列的向正面迅速强攻，因其从我侧翼迂回又必须先向翼侧筑垒，这样不唯消耗兵力，而且消耗时间，对于我们终是有利的；（二）不愿远离堡垒施行追击，即战场追击亦很少，只要夺取了某个目标后，立即筑垒修路，这样使每进攻一次须得停顿三天、五天，有时半月，因为如此便给了我们在纵深地带上有新的筑阵地的可能，这次战斗敌人攻占不到三个启罗米达的水平距离，共费三天的时间（二次跃进与两天剧战），炮弹约二千发，投弹三百枚，小枪小弹总数十倍于我，死伤兵员约三千，只团长一级的上级军官就有五个，但我们仍在他的前面构筑了新的阵地，假如我们在几百里距离的赤色版图上，一开始就使敌人遭受这样的抵抗，而给敌人消耗量当是不可计算的，要记着广昌战斗我们有生力量的消耗是数倍于敌的。

守备部队向来是我们用较脆弱的兵团来担任的，而精干的兵团都用于突击的方向，因此将乐、泰宁的失守，东华山的失守，广昌及广昌以北地区诸要点的失守，大寨脑的失守，都使我们突击未成而支点已失，结果突击计划形成泡影，有时可能遭受意外的打击（如东华山）……但太阳嶂、马鞍寨、高虎脑，因为守备队是从基干兵团内发出来的，虽然兵力是不很多，然而都充分显示了他的坚强防御，所以在扼守某些决定胜负的地段，分配兵力的任务关系上，应该得到教训，高脚岭与高虎脑也是这个战斗中的具体例子。\footnote{《三军团半桥附近战斗详报（1934年8月5日至7日）》，《中央苏区第五次反“围剿”》（下），第102～103页。}
\end{quote}


客观地看，彭德怀的这一意见并不是毫无道理，当然，如果真按这样一种思路去进行反“围剿”战争，也未必会有好的结果，它只是再一次提醒我们，历史是如此复杂，在历史的此和彼之间，也许还有着更多的曲径幽道。

